\chapter{让一次失败告诉我们下一次怎么选}\label{ch:descent-finish}
\section{分类之前，先看我们在怕什么}
现在有两个不同重值 $a,b$。如果完整留住 $p$ 份 $b$，另找一组做整数平均，唯一担心是新均值又等于 $b$。如果它不等于 $b$，我们就有新旧两个重值，任务完成。

这看似是麻烦，其实是额外信息。设 $a=b+\delta$，尝试
\[(a^i,b^{p-1-i},x).
\]
总和为 $pb+i\delta+x-b$。当 $p\nmid\delta$ 时，模 $p$ 可以唯一选择 $i\in\{0,\ldots,p-1\}$ 使平均为整数。如果均值恰为 $b$，就不只是一条同余，而是严格等式
\[x=b-i\delta.
\]
所以“所有尝试都失败”会强迫所有第三值落在一段有限等差列里。这时零和条件往往就不容许它们继续失败。

\section{副本多时：先拿背景值补足人数}
如果 $b$ 有至少 $2p-1$ 份，我们最多拿走 $p-1$ 份，仍留下一个完整 $p$ 块。因此它是很好的“补足人数”的背景值。

下面证明里出现的长选择式有一个简单来源。负方向的值是 $b-j_t\delta$，系数总和第一次达到 $T\ge p$。我们希望新均值为 $b-\delta$，即所选的相对系数和为 $-p$。已经拿了系数和 $-T$，就补 $T-p$ 份系数为 $1$ 的 $a$；余下位置用系数为零的 $b$ 填满。于是得到
\[(b-j_1\delta,\ldots,b-j_s\delta,a^{T-p},b^{2p-s-T}).
\]
不是凭空猜到这个式子，而是“先指定总和，再补足项数”。证明中的 $T+s\le2p$ 正好保证最后的补足数不为负。
\InsertProof{lem:background}

\section{重值种类多时：数的大小顺序就够用了}
如果有四个不同重值，先把最大的一个留在外面，只用另外三个。无论怎样平均，新均值都严格小于那个最大值，碰撞自动被排除。

只剩三个重值时，位置稍紧，需要先看有没有轻值、极端值有没有多出一份。最后的困难形状是两端各恰有 $p$ 份，中间有 $p+r$ 份。先拿一份中间值；如果碰撞，再拿两份。两次碰撞方程不能同时成立，除非已经是相反数加零的终端。
\InsertProof{lem:heavy}
这个引理后半还处理“两个重值模 $p$ 相等”的情况。它用的部分和碰撞是抽屉原理：$p+1$ 个部分和只有 $p$ 种余数，总有两个相同；它们之间那一段就是一个非空零和片段。

\section{两边各多一点，也可以合起来用}
不必总要求某一个重值单独有 $2p-1$ 份。设它们分别有 $p+\alpha,p+\beta$ 份，我们把 $\alpha,\beta$ 叫作富余份数。如果一次总共只从两边拿 $p-2$ 份，而 $\alpha+\beta\ge p-2$，就不可能同时把两边都拿到不足 $p$ 份。

这个简单计数允许两个重值共同承担背景值的工作。下面的等差结构依然由“假设所有单轻值选择都碰撞”推出来，再由零和排除。
\InsertProof{lem:two-heavy}

\section{富余也不够时：借出一份，再借出两份}
最后只剩两个不同余重值，并且富余总数最多 $p-3$。此时先各留 $p$ 份，再保护最多 $d_p(n)$ 个位置。如果
\[n\ge4p-2+d_p(n),
\]
还能拿出 $2p-2$ 个候选。上一讲说，若找不到 $p$ 项零和，这个池只能分成两类，各 $p-1$ 份。

现在借一份锚点 $a$ 进来，用两类余数补足。若所有实际选择都碰撞到 $b$，在同一余数类内交换“拿了谁、没拿谁”，碰撞等式会强迫同类位置的实际数值都相等。接着借两份 $a$，便能打破碰撞。若唯一的对称选择不允许借两份，就改从 $b$ 端做；两端不可能同时处于那个坏位置。
\InsertProof{lem:integer-range}

这里终于可以使用能量终止：已经证明每个非终端都有下一步，而且每步都使非负整数 $E$ 严格减小。因此过程不会无路可走，也不能无限走下去。这两件事都要有，单独一句“能量下降”还不够。

我们得到的是一块宽维数范围。它已经很大，但尚未达到 $2p+1$。后面会看到，位置紧张时该换一种办法：允许分数，用精确输送代替每步整数下降。
